% Options for packages loaded elsewhere
\PassOptionsToPackage{unicode}{hyperref}
\PassOptionsToPackage{hyphens}{url}
\documentclass[
  spanish,
  12pt,
  letterpaper,
]{article}
\usepackage{xcolor}
\usepackage[margin=2.4cm]{geometry}
\usepackage{amsmath,amssymb}
\setcounter{secnumdepth}{-\maxdimen} % remove section numbering
\usepackage{iftex}
\ifPDFTeX
  \usepackage[T1]{fontenc}
  \usepackage[utf8]{inputenc}
  \usepackage{textcomp} % provide euro and other symbols
\else % if luatex or xetex
  \usepackage{unicode-math} % this also loads fontspec
  \defaultfontfeatures{Scale=MatchLowercase}
  \defaultfontfeatures[\rmfamily]{Ligatures=TeX,Scale=1}
\fi
\usepackage{lmodern}
\ifPDFTeX\else
  % xetex/luatex font selection
  \setmainfont[]{Times New Roman}
\fi
% Use upquote if available, for straight quotes in verbatim environments
\IfFileExists{upquote.sty}{\usepackage{upquote}}{}
\IfFileExists{microtype.sty}{% use microtype if available
  \usepackage[]{microtype}
  \UseMicrotypeSet[protrusion]{basicmath} % disable protrusion for tt fonts
}{}
\usepackage{setspace}
\makeatletter
\@ifundefined{KOMAClassName}{% if non-KOMA class
  \IfFileExists{parskip.sty}{%
    \usepackage{parskip}
  }{% else
    \setlength{\parindent}{0pt}
    \setlength{\parskip}{6pt plus 2pt minus 1pt}}
}{% if KOMA class
  \KOMAoptions{parskip=half}}
\makeatother
\ifLuaTeX
\usepackage[bidi=basic,shorthands=off,]{babel}
\else
\usepackage[bidi=default,shorthands=off,]{babel}
\fi
\ifPDFTeX
\else
\babelfont{rm}[]{Times New Roman}
\fi
\ifLuaTeX
  \usepackage{selnolig} % disable illegal ligatures
\fi
\setlength{\emergencystretch}{3em} % prevent overfull lines
\providecommand{\tightlist}{%
  \setlength{\itemsep}{0pt}\setlength{\parskip}{0pt}}

\usepackage{titlesec}
\titleformat{\section}{\centering\bfseries}{\thesection}{1em}{}
\titleformat{\subsection}{\raggedright\bfseries}{\thesubsection}{1em}{}
\titleformat{\subsubsection}{\raggedright\bfseries\itshape}{\thesubsubsection}{1em}{}
\titleformat{\paragraph}[runin]{\raggedright\bfseries}{\theparagraph}{1em}{}[.]
\titleformat{\subparagraph}[runin]{\raggedright\bfseries\itshape}{\thesubparagraph}{1em}{}
\titlespacing*{\paragraph}{0cm}{*1}{1em}
\titlespacing*{\subparagraph}{0cm}{*1}{1em}

\setlength{\parindent}{0cm}
\setlength{\parskip}{0pt}

\usepackage{caption}
\captionsetup[table]{labelfont=bf,textfont=it,labelsep=newline,justification=raggedright,singlelinecheck=false}
\captionsetup[figure]{labelfont=bf,textfont=it,labelsep=newline,justification=raggedright,singlelinecheck=false}

\usepackage{xurl}
\tolerance=9999
\emergencystretch=3em
\hbadness=10000

\usepackage{bookmark}
\IfFileExists{xurl.sty}{\usepackage{xurl}}{} % add URL line breaks if available
\urlstyle{same}
\hypersetup{
  pdflang={es},
  hidelinks,
  pdfcreator={LaTeX via pandoc}}

\author{}
\date{}

\begin{document}

\setstretch{1}
\begin{center}

\textbf{SOFTWARE Y LAS ARQUITECTURAS}

\emph{D. A. Flores Cordero}

7690-14-745 Universidad Mariano Gálvez

047 Seminario de Tecnologías de Información

\emph{dfloresc3@miumg.edu.gt}

\end{center}

\section{Resumen}\label{resumen}

Este trabajo investiga las ``arquitecturas tipo'' documentadas para
aplicaciones de software, con el objetivo de identificar qué criterio
real distingue a un monolito, un modelo N-capas, una arquitectura limpia
(hexagonal) y una arquitectura de microservicios, más allá de sus
nombres. La investigación partió de tres fuentes: la guía de
arquitecturas de aplicaciones web de Microsoft, que documenta la
evolución de un proyecto único hacia capas, arquitectura limpia y
contenedores; una propuesta académica de modelo en cinco capas para
aplicaciones web multiplataforma; y un marco de evaluación de
arquitectura de aplicaciones centrado en cuatro elementos
---funcionalidad, ensamblaje, interacción y gestión de datos---. A
partir de estas fuentes se identificó que las arquitecturas tipo no son
alternativas que compiten entre sí, sino puntos distintos en una misma
escala de separación de responsabilidades, y que su costo real no está
en el nombre del patrón sino en las dependencias que genera entre capas
y en la coordinación que exige entre servicios. Los resultados muestran
que la arquitectura limpia resuelve el problema de dependencia
direccional del modelo N-capas tradicional, y que los microservicios
solo se justifican cuando existe una necesidad real de escalar
componentes de forma independiente, no como estándar por defecto. Se
concluye que elegir una arquitectura tipo es una decisión de
costo-beneficio ligada al tamaño y a la concurrencia real del sistema,
no una preferencia estética entre patrones.

\textbf{Palabras claves:} arquitectura de software, monolito,
arquitectura por capas, arquitectura limpia, microservicios

\section{Desarrollo del tema}\label{desarrollo-del-tema}

Toda aplicación profesional necesita un límite entre sus
responsabilidades; la única decisión real es en qué momento y con qué
costo se traza ese límite. Las arquitecturas documentadas en la
literatura técnica ---monolito de un solo proyecto, N-capas,
arquitectura limpia (u hexagonal) y microservicios--- no compiten por
ser ``la correcta'': son puntos distintos en la misma escala de
separación de responsabilidades, y elegir la que no corresponde al
tamaño real del problema sale más caro que no elegir ninguna (Microsoft,
s.f.).

\subsection{El monolito como punto de
partida}\label{el-monolito-como-punto-de-partida}

Un proyecto ASP.NET Core nuevo comienza, por defecto, como un monolito
``todo en uno'': presentación, lógica de negocio y acceso a datos
conviven en un mismo ensamblado, separados apenas por carpetas
(Microsoft, s.f.). Esta simplicidad tiene un límite temporal, no
estructural: a medida que crece el número de archivos, la falta de una
frontera real entre carpetas ---y no solo entre nombres de carpetas---
deriva en código espagueti, porque nada impide que una vista llame
directamente a la capa de datos. La respuesta habitual no es abandonar
el monolito, sino dividirlo en varios proyectos por responsabilidad
dentro del mismo despliegue. Vale además una distinción que rara vez se
hace explícita: capas es separación lógica dentro del código; niveles es
distribución física en servidores o procesos distintos. Es perfectamente
posible ---y común--- tener una aplicación de varias capas que se
implementa en un único nivel (Microsoft, s.f.).

\subsection{El modelo N-capas y su costo de
dependencia}\label{el-modelo-n-capas-y-su-costo-de-dependencia}

La organización más citada divide la aplicación en interfaz de usuario
(UI), lógica de negocio (BLL) y acceso a datos (DAL), donde cada capa
solo puede llamar a la inmediata inferior. El problema que rara vez se
menciona junto a este diagrama es que las dependencias de compilación
corren de arriba hacia abajo: la lógica de negocio, que debería ser el
corazón de la aplicación, termina dependiendo de los detalles de acceso
a datos ---y, con frecuencia, de la existencia de una base de datos real
para poder probarla (Microsoft, s.f.). Una propuesta académica
alternativa lleva esta separación un paso más allá con un modelo de
cinco capas explícitas ---base de datos, lógica de manejo de datos,
procesos de negocio, control de interfaz y vista---, donde cada capa se
comunica con la siguiente mediante formatos neutros (XML, JSON) sobre
HTTP, de modo que cualquier capa pueda sustituirse sin reescribir las
demás (Gómez Fermín \& Moreno Poggio, 2014). El modelo admite distintos
despliegues: monolítico en un solo servidor, escalabilidad lineal con
una capa por servidor, escalabilidad independiente con réplicas por
capa, o clúster para la base de datos. La ganancia en flexibilidad y
disponibilidad es real, pero también lo es el costo: más servidores, más
piezas que administrar y, en el caso de la escalabilidad independiente,
un único punto de falla si la capa de base de datos cae, sin importar
cuántas réplicas tengan las capas superiores (Gómez Fermín \& Moreno
Poggio, 2014).

\subsection{Arquitectura limpia: invertir la
dependencia}\label{arquitectura-limpia-invertir-la-dependencia}

La arquitectura limpia ---también conocida como hexagonal, de puertos y
adaptadores, u onion--- resuelve exactamente el problema de dirección
que arrastra el modelo N-capas tradicional. En lugar de que la lógica de
negocio dependa del acceso a datos, esa dependencia se invierte: el
núcleo de la aplicación define interfaces, y es la infraestructura la
que depende de ese núcleo al implementarlas (Microsoft, s.f.).
Visualizada como círculos concéntricos, en el centro están las entidades
y los servicios de dominio; afuera, interfaz de usuario e
infraestructura dependen del núcleo, pero no dependen entre sí. La
ventaja concreta no es solo conceptual: al no depender de
infraestructura, el núcleo de la aplicación se puede probar con pruebas
unitarias sin necesidad de una base de datos real, y las
implementaciones de infraestructura se pueden intercambiar ---de SQL
Server a una API en la nube, por ejemplo--- sin tocar la lógica de
negocio (Microsoft, s.f.). Sin embargo, esta inversión de dependencias
no es gratuita: exige definir interfaces y una raíz de composición donde
se conectan las implementaciones concretas, una capa adicional de
indirección que un equipo pequeño puede no necesitar para una aplicación
simple.

\subsection{Contenedores y microservicios: independencia con costo de
coordinación}\label{contenedores-y-microservicios-independencia-con-costo-de-coordinaciuxf3n}

Un monolito también puede desplegarse como un único contenedor Docker,
lo cual resuelve el problema de ``funciona en mi máquina'' y agiliza la
actualización mediante imágenes inmutables que se levantan en segundos
(Microsoft, s.f.). El límite aparece con el escalado: en una tienda en
línea, muchos más clientes navegan el catálogo de productos que los que
llegan a pagar, pero un monolito escala todo el código por igual, aunque
solo un componente sea el cuello de botella. Ahí es donde los
microservicios ofrecen una ventaja real: cada servicio se escala, se
prueba y se despliega de forma independiente, dentro de un contexto
delimitado propio. La contraparte es que separar una aplicación en
procesos discretos agrega complejidad de comunicación ---ya no hay
llamadas a métodos, sino mensajería asíncrona--- y obliga a resolver
control de bus de eventos, reintentos, resiliencia de mensajes y
consistencia eventual, piezas que un monolito nunca necesitó (Microsoft,
s.f.). Esta misma lógica de desacoplar componentes mediante un
middleware neutro a la plataforma ---sin que un servicio conozca la
implementación interna del otro--- es la que sostiene también al modelo
de cinco capas descrito antes (Gómez Fermín \& Moreno Poggio, 2014).

\subsection{Un marco para evaluar, no solo para
nombrar}\label{un-marco-para-evaluar-no-solo-para-nombrar}

Más allá del nombre del patrón, cualquier arquitectura de aplicación se
puede evaluar con cuatro elementos: funcionalidad (qué hace el sistema),
ensamblaje (cómo se agrupan sus partes), interacción (cómo se comunican)
y gestión de datos (cómo se almacenan, protegen y mueven los datos)
(Conexiam, s.f.). El ensamblaje es, según esta perspectiva, el desafío
crítico: es donde se decide si la funcionalidad vive integrada en una
sola aplicación, expuesta como microservicio, o distribuida en varias
capas, y esa decisión de ensamblaje es la que define los límites de
integración y el ciclo de vida de cada pieza (Conexiam, s.f.). Este
marco tiene una utilidad práctica que los diagramas de capas no ofrecen
por sí solos: permite comparar un monolito con una arquitectura limpia o
con microservicios usando los mismos cuatro criterios, en lugar de
discutir cuál ``se ve'' más moderna.

\section{Observaciones y comentarios}\label{observaciones-y-comentarios}

Investigar tres fuentes con enfoques distintos ---una guía de
implementación de Microsoft, una propuesta académica de modelo en capas
y un marco de evaluación de arquitectura empresarial--- confirmó que la
mayoría de la documentación técnica describe cómo estructurar el código,
pero rara vez cuantifica el costo operativo de esa estructura. El
hallazgo más útil no fue una arquitectura específica, sino el marco de
cuatro elementos de Conexiam (s.f.): funciona como vocabulario común
para comparar patrones que, a simple vista, parecen inconmensurables
entre sí.

\section{Conclusiones}\label{conclusiones}

\begin{enumerate}
\def\labelenumi{\arabic{enumi}.}
\tightlist
\item
  Monolito, N-capas, arquitectura limpia y microservicios no son
  alternativas excluyentes: son puntos distintos en la misma escala de
  separación de responsabilidades.
\item
  Capas es separación lógica del código; niveles es distribución física
  en servidores. Una aplicación de varias capas puede desplegarse
  perfectamente en un solo nivel.
\item
  La arquitectura limpia resuelve el problema de dependencia direccional
  del modelo N-capas tradicional invirtiendo las dependencias hacia el
  núcleo de negocio, a costa de una capa adicional de indirección.
\item
  Los microservicios solo se justifican cuando existe una necesidad real
  de escalar componentes de forma independiente; de lo contrario, la
  complejidad de coordinación que agregan supera el beneficio.
\item
  Evaluar arquitecturas por funcionalidad, ensamblaje, interacción y
  gestión de datos es más útil que compararlas únicamente por su nombre
  o popularidad.
\end{enumerate}

\section{Bibliografía}\label{bibliografuxeda}

Conexiam. (s.f.). \emph{¿Qué es la arquitectura de aplicaciones?}.
https://conexiam.com/es/que-es-la-arquitectura-de-aplicaciones/

Gómez Fermín, L. V., \& Moreno Poggio, T. R. (2014). \emph{Propuesta de
modelo en cinco capas para aplicaciones web}. SciELO Venezuela.
https://ve.scielo.org/scielo.php?script=sci\_arttext\&pid=S1315-01622014000200009

Microsoft. (s.f.). \emph{Arquitecturas de aplicaciones web comunes}.
Microsoft Learn.
https://learn.microsoft.com/es-es/dotnet/architecture/modern-web-apps-azure/common-web-application-architectures

\begin{center}\rule{0.5\linewidth}{0.5pt}\end{center}

Repositorio del documento (código fuente LaTeX):
\url{https://github.com/dan-desa/seminario-tecnologias-2026/tree/main/foro-03}

\end{document}
